\documentclass[12pt]{article}
\usepackage[T1]{fontenc}
\usepackage[utf8]{inputenc}
\usepackage{lmodern}
\usepackage{textcomp}
\usepackage{enumitem}
\usepackage{geometry}
\geometry{margin=1in}
\begin{document}
\begin{center}
Answer Key - Psychology - Graduate Level Assessment 1 \
\small (Answers are ordered exactly as in the question paper.)
\end{center}

\section*{Multiple Choice}
\begin{enumerate}[label=\arabic*.]
\item \textbf{Answer:} C
\\ \textit{Options:} A) reciprocal determinism; B) cognitive dynamogenesis; C) psychic determinism; D) ptyalism; E) linguistic determinism
\\ \textit{Note:} The correct answer is psychic determinism because it posits that all mental processes, including dreams and slips of the tongue, are driven by underlying psychological motivations and serve specific goals, reflecting Freud's belief that no behavior is random or without purpose.
\item \textbf{Answer:} A
\\ \textit{Options:} A) tend to prefer a less directive therapeutic approach.; B) are grounded more in the here-and-now than in the past or the future.; C) are likely to express emotional problems as verbal symptoms.; D) respond better when goal-setting is immediate.; E) are grounded more in the past or the future than in the here-and-now.
\\ \textit{Note:} Research indicates that Asian and Asian-American clients often value harmony and indirect communication, leading them to prefer a less directive therapeutic approach that allows for more subtle exploration of issues.
\item \textbf{Answer:} B
\\ \textit{Options:} A) Adlerian; B) reality; C) person-centered; D) existential; E) Gestalt
\\ \textit{Note:} Reality therapy focuses on helping individuals recognize and change their self-perceptions and behaviors by replacing negative self-identities, such as a "failure identity," with positive and empowering self-identities, such as a "success identity."
\item \textbf{Answer:} D
\\ \textit{Options:} A) CAT; B) X-ray; C) EEG; D) PET; E) MRI
\\ \textit{Note:} PET scans measure metabolic activity and brain function, providing evidence that significant regions of the brain are active during various tasks, thereby disproving the myth that only 10 percent is utilized.
\item \textbf{Answer:} D
\\ \textit{Options:} A) not mentioned in the ethical guidelines.; B) ethically acceptable, as it provides opportunities for those with lower income.; C) explicitly recommended in the ethical guidelines.; D) ethically acceptable since it serves the best interests of his clients.; E) explicitly recommended in the legal guidelines.
\\ \textit{Note:} Dr. Manzetti's use of a sliding scale based on clients' income demonstrates ethical practice by ensuring that therapy remains accessible and accommodating to clients' financial situations, ultimately prioritizing their best interests and well-being.
\end{enumerate}

\section*{True / False}
\begin{enumerate}[label=\arabic*.]
\setcounter{enumi}{5}
\item \textbf{Answer:} False
\\ \textit{Options:} A) True; B) False
\\ \textit{Note:} The statement is false because the mean weight of a sample does not necessarily represent the average weight of the entire population of rats, as it is calculated from only a subset of the population, which may not be representative.
\item \textbf{Answer:} False
\\ \textit{Options:} A) True; B) False
\\ \textit{Note:} Psychology primarily studies individual behavior and mental processes rather than focusing on the historical analysis of human events and their societal impact.
\item \textbf{Answer:} True
\\ \textit{Options:} A) True; B) False
\\ \textit{Note:} Vestibule training is effective for jobs like receptionist because it allows individuals to practice and refine their skills in a controlled, realistic setting, enhancing their readiness for actual job responsibilities.
\item \textbf{Answer:} True
\\ \textit{Options:} A) True; B) False
\\ \textit{Note:} The statement is true because serving as both a fact witness and an expert witness creates a conflict of interest that can compromise the objectivity and integrity of the psychologist's testimony, violating ethical guidelines that require clear boundaries between these roles.
\item \textbf{Answer:} True
\\ \textit{Options:} A) True; B) False
\\ \textit{Note:} The statement is true because a mode is defined as the number that appears most frequently in a data set, and in this sample, all numbers appear with the same frequency of one occurrence each, resulting in no mode.
\end{enumerate}

\section*{Long Form Response}
\begin{enumerate}[label=\arabic*.]
\setcounter{enumi}{10}
\item \textbf{Answer:} Carl Rogers' perspective on personality centers on the belief that individuals possess an inherent positivity and a goal-directed nature, which he termed the actualizing tendency. This concept posits that people are fundamentally oriented towards growth, self- fulfillment, and the realization of their potential, suggesting that personality traits are not merely reactive but actively pursue constructive outcomes. The implications for psychological growth are profound, as Rogers' approach encourages individuals to embrace their innate capacities for change and development, fostering a therapeutic environment rooted in acceptance and empathy. In therapy, this perspective facilitates a client- centered approach, where the therapist supports clients in uncovering their strengths and aligning their behaviors with their authentic self, ultimately promoting healthier personality development and psychological well-being.
\\ \textit{Note:} correct because it accurately captures Carl Rogers' concept of the actualizing tendency, emphasizing that personality traits are inherently positive and goal-directed, which significantly influences psychological growth and therapeutic practices by promoting self-fulfillment and potential realization.
\item \textbf{Answer:} The classic aging pattern in cognitive assessment, particularly as observed in the performance subtests of the Wechsler Adult Intelligence Scale-Revised (WAIS-R), reveals a decline in specific cognitive domains such as processing speed, perceptual organization, and fluid intelligence. Older adults typically demonstrate preserved verbal abilities while experiencing more significant deficits in performance subtests, which assess skills like block design and digit symbol coding. This differential decline suggests that while crystallized intelligence may remain stable, the more novel and adaptive aspects of cognition tend to diminish with age. Understanding these patterns is crucial for clinicians and researchers as they navigate the complexities of cognitive aging, providing insight into the specific challenges older adults face and informing interventions aimed at maintaining cognitive health.
\\ \textit{Note:} The answer correctly identifies the classic aging pattern in cognitive assessment by emphasizing the decline in performance subtests related to processing speed and fluid intelligence in older adults, while noting the preservation of verbal abilities, thereby illustrating the differential impact of aging on various cognitive domains.
\end{enumerate}

\vfill
\noindent\textit{End of Answer Key}
\end{document}